% !TEX program = xelatex
% ============================================================
%  اتفاقية المساهمين — Shareholder Agreement Template
%  نموذج اتفاقية مساهمين تنظم العلاقة بين حملة الأسهم في الشركة
%  Compile with: xelatex main.tex (run twice)
% ============================================================
%  Features:
%    - أطراف ومقدمات (Parties & Recitals)
%    - مواد مرقمة: رأس المال، الحقوق، نقل الأسهم، حقوق الأولوية، مجلس الإدارة، الاجتماعات، الخروج
%    - جدول رأس المال وتوزيع الأسهم (Share Capital Table)
%    - جدول حقوق الأولوية (Pre-emptive Rights Table)
%    - كتلة توقيعات للمساهمين (Signatures Block)
%    - نص عربي قانوني مع مصطلحات إنجليزية بين قوسين
% ============================================================

\documentclass[12pt,a4paper]{article}

\usepackage[a4paper,margin=2.5cm,top=2.5cm,bottom=2.5cm,headheight=16pt]{geometry}
\usepackage{graphicx}
\usepackage{array}
\usepackage{booktabs}
\usepackage{tabularx}
\usepackage{multirow}
\usepackage{enumitem}
\usepackage{float}
\usepackage{titlesec}
\usepackage{fancyhdr}
\usepackage{setspace}
\usepackage{xcolor}

\usepackage{bidi}
\usepackage[no-math]{fontspec}
\usepackage[quiet]{polyglossia}

\setmainfont[Scale=1]{NotoKufiArabic-Regular.ttf}
\newfontfamily\arabicfont[Script=Arabic,Scale=0.95]{NotoKufiArabic-Regular.ttf}
\newfontfamily\englishfont[Scale=0.95]{TeX Gyre Termes}

\setdefaultlanguage[calendar=gregorian,hijricorrection=1]{arabic}
\setotherlanguage[variant=british]{english}

\usepackage[breaklinks,hidelinks]{hyperref}
\hypersetup{pdftitle={اتفاقية المساهمين},pdfauthor={اسم الشركة}}

\titleformat{\section}{\large\bfseries}{\thesection}{0.7em}{}
\titlespacing*{\section}{0pt}{1.2ex plus 0.3ex}{0.8ex plus 0.2ex}

\pagestyle{fancy}
\fancyhf{}
\fancyhead[R]{\small\itshape اتفاقية المساهمين (Shareholder Agreement)}
\fancyhead[L]{\small اسم الشركة}
\fancyfoot[C]{\small\thepage}
\renewcommand{\headrulewidth}{0.3pt}

\newcolumntype{L}[1]{>{\raggedright\arraybackslash}p{#1}}
\newcolumntype{C}[1]{>{\centering\arraybackslash}p{#1}}

\setstretch{1.3}

\begin{document}

\begin{center}
{\LARGE\bfseries اتفاقية المساهمين}\\[0.5em]
{\large (\LR{Shareholder Agreement})}\\[1em]
\end{center}

\vspace{0.5em}

% TODO: استبدل التاريخ والمرجع بالقيم الفعلية
\noindent\textbf{رقم الاتفاقية:} \rule{4cm}{0.4pt} \hfill \textbf{التاريخ:} \rule{4cm}{0.4pt}

\vspace{1em}

% ============================================================
% PARTIES
% ============================================================
\noindent\textbf{الأطراف (\LR{Parties}):}

\noindent وقّع المساهمون أدناه على هذه الاتفاقية، ويُشار إليهم مجتمعين بـ«المساهمين» ومنفردين بـ«المساهم»:

\begin{enumerate}
    \item \textbf{المساهم الأول (\LR{Shareholder 1}):} \rule{6cm}{0.4pt}، رقم الهوية/السجل التجاري: \rule{4cm}{0.4pt}، العنوان: \rule{5cm}{0.4pt}.

    \item \textbf{المساهم الثاني (\LR{Shareholder 2}):} \rule{6cm}{0.4pt}، رقم الهوية/السجل التجاري: \rule{4cm}{0.4pt}، العنوان: \rule{5cm}{0.4pt}.

    % TODO: أضف مساهمين إضافيين حسب الحاجة
\end{enumerate}

\vspace{1em}

% ============================================================
% RECITALS
% ============================================================
\noindent\textbf{الديباجة (\LR{Recitals}):}

\begin{enumerate}
    \item الشركة المسجّلة بموجب السجل التجاري رقم \rule{4cm}{0.4pt} باسم \rule{6cm}{0.4pt} (يُشار إليها بـ«الشركة»).
    \item يملك المساهمون أسهماً في الشركة وفقاً للجدول الوارد في المادة الأولى.
    \item يرغب المساهمون في تنظيم علاقتهم وإدارة الشركة وفقاً لشروط هذه الاتفاقية.
\end{enumerate}

\vspace{1em}

% ============================================================
% 1. SHARE CAPITAL
% ============================================================
\section{المادة الأولى: رأس المال وتوزيع الأسهم (\LR{Share Capital})}
\label{sec:capital}
يبلغ رأس المال المصرّح به \rule{3cm}{0.4pt} سهماً، ورأس المال المُصدر \rule{3cm}{0.4pt} سهماً، موزّعة كما يلي:

% TODO: املأ تفاصيل توزيع الأسهم
\begin{table}[H]
\centering
\renewcommand{\arraystretch}{1.5}
\begin{tabularx}{\textwidth}{L{5cm} C{3cm} C{3cm} C{2.5cm}}
\toprule
\textbf{المساهم} & \textbf{عدد الأسهم} & \textbf{القيمة الاسمية} & \textbf{النسبة} \\
\midrule
المساهم الأول & \rule{2cm}{0.4pt} & \rule{2cm}{0.4pt} & \rule{1.5cm}{0.4pt}\% \\
المساهم الثاني & \rule{2cm}{0.4pt} & \rule{2cm}{0.4pt} & \rule{1.5cm}{0.4pt}\% \\
\midrule
\textbf{الإجمالي} & \textbf{\rule{2cm}{0.4pt}} & \textbf{---} & \textbf{100\%} \\
\bottomrule
\end{tabularx}
\end{table}

% ============================================================
% 2. RIGHTS
% ============================================================
\section{المادة الثانية: حقوق المساهمين (\LR{Shareholder Rights})}
\label{sec:rights}
يتمتع كل مساهم بالحقوق التالية:
\begin{enumerate}
    \item \textbf{حق التصويت (\LR{Voting Rights}):} لكل سهم صوت واحد في الجمعية العامة.
    \item \textbf{حق توزيعات الأرباح (\LR{Dividend Rights}):} الحصول على حصته من الأرباح الموزّعة.
    \item \textbf{حق الاطلاع (\LR{Information Rights}):} الاطلاع على القوائم المالية وسجلات الشركة.
    \item \textbf{حق نقل الأسهم (\LR{Transfer Rights}):} نقل أسهمه وفقاً لشروط المادة الثالثة.
    \item \textbf{حقوق الأقلية (\LR{Minority Rights}):} الحماية وفقاً للقوانين السارية.
\end{enumerate}

% ============================================================
% 3. TRANSFER OF SHARES
% ============================================================
\section{المادة الثالثة: نقل الأسهم (\LR{Transfer of Shares})}
\label{sec:transfer}
\begin{enumerate}
    \item لا يجوز لأي مساهم نقل أسهمه (\LR{Transfer Shares}) إلى طرف ثالث دون التزام بإجراءات هذه المادة.
    \item يجب على المساهم الراغب في النقل إخطار الشركة والمساهمين الآخرين كتابةً بعرض البيع (\LR{Notice of Sale}).
    \item يتمتع المساهمون الآخرون بحق الأولوية (\LR{Right of First Refusal}) لشراء الأسهم بالسعر والشروط نفسها.
    \item في حال عدم ممارسة حق الأولوية خلال \rule{2cm}{0.4pt} يوماً، يجوز النقل للطرف الثالث.
    \item يخضع أي نقل لموافقة الجمعية العامة وفقاً للقوانين السارية.
\end{enumerate}

% ============================================================
% 4. PRE-EMPTIVE RIGHTS
% ============================================================
\section{المادة الرابعة: حقوق الأولوية في الإصدارات الجديدة (\LR{Pre-emptive Rights})}
\label{sec:preemptive}
\begin{enumerate}
    \item عند إصدار الشركة أسهماً جديدة (\LR{New Issue})، يتمتع كل مساهم بحق الأولوية في الاكتتاب بنسبة حصته الحالية.
    \item يجب على المساهم ممارسة هذا الحق خلال \rule{2cm}{0.4pt} يوماً من تاريخ الإخطار.

% TODO: حدّد شروط الإصدارات الجديدة
\begin{table}[H]
\centering
\renewcommand{\arraystretch}{1.5}
\begin{tabularx}{\textwidth}{L{6cm} C{3cm} C{3cm}}
\toprule
\textbf{البند} & \textbf{المساهم الأول} & \textbf{المساهم الثاني} \\
\midrule
نسبة الأولوية في الإصدار الجديد & \rule{2cm}{0.4pt}\% & \rule{2cm}{0.4pt}\% \\
مدة ممارسة الحق (يوم) & \rule{2cm}{0.4pt} & \rule{2cm}{0.4pt} \\
\bottomrule
\end{tabularx}
\end{table}

    \item في حال عدم ممارسة الحق، يجوز للشركة طرح الأسهم على أطراف ثالثة.
\end{enumerate}

% ============================================================
% 5. BOARD OF DIRECTORS
% ============================================================
\section{المادة الخامسة: مجلس الإدارة (\LR{Board of Directors})}
\label{sec:board}
\begin{enumerate}
    \item يُكوَّن مجلس الإدارة من \rule{2cm}{0.4pt} أعضاء يُنتخبون من قبل المساهمين.
    \item يحق لكل مساهم ترشيح عدد من المرشحين يتناسب مع نسبة ملكيته.
    \item تُتخذ قرارات المجلس بالأغلبية، وتتطلب القرارات الكبرى أغلبية مؤهلة (\LR{Supermajority}) بنسبة \rule{2cm}{0.4pt}\%.
    \item يُعقد المجلس اجتماعات دورية \rule{2cm}{0.4pt} (شهرياً/ربعياً) ويدوّن محاضرها.
\end{enumerate}

% ============================================================
% 6. MEETINGS
% ============================================================
\section{المادة السادسة: اجتماعات المساهمين (\LR{Shareholder Meetings})}
\label{sec:meetings}
\begin{enumerate}
    \item \textbf{الجمعية العامة السنوية (\LR{Annual General Meeting}):} تُعقد مرة سنوياً خلال \rule{2cm}{0.4pt} أشهر من نهاية السنة المالية.
    \item \textbf{الجمعية العامة غير العادية (\LR{Extraordinary General Meeting}):} تُعقد عند الحاجة بطلب من \rule{3cm}{0.4pt}\% من المساهمين.
    \item \textbf{النصاب القانوني (\LR{Quorum}):} يُشترط حضور \rule{2cm}{0.4pt}\% من رأس المال لانعقاد الاجتماع.
    \item تُرسل الدعوات للاجتماعات قبل \rule{2cm}{0.4pt} يوماً من تاريخ الانعقاد.
\end{enumerate}

% ============================================================
% 7. EXIT
% ============================================================
\section{المادة السابعة: حقوق الخروج (\LR{Exit Rights})}
\label{sec:exit}
\begin{enumerate}
    \item \textbf{حق المرافقة (\LR{Tag-along Right}):} إذا رغب مساهم يملك أغلبية الأسهم في البيع، يحق للأقلية مرافقته في البيع بالشروط نفسها.
    \item \textbf{حق السحب (\LR{Drag-along Right}):} إذا وافق \rule{2cm}{0.4pt}\% من المساهمين على البيع، يُلزم الباقي بالبيع بالشروط نفسها.
    \item \textbf{الاكتتاب العام (\LR{IPO}):} يجوز للشركة طرح أسهمها للاكتتاب العام بقرار من الجمعية العامة.
    \item \textbf{التصفية (\LR{Liquidation}):} عند تصفية الشركة، تُوزَّع الأصول المتبقية بعد سداد الالتزامات بنسب الملكية.
\end{enumerate}

% ============================================================
% SIGNATURES
% ============================================================
\vspace{2em}
\section*{التوقيعات (\LR{Signatures})}

% TODO: املأ بيانات الموقّعين
\begin{table}[H]
\centering
\renewcommand{\arraystretch}{2.5}
\begin{tabularx}{\textwidth}{L{5cm} X X}
\toprule
& \textbf{المساهم الأول} & \textbf{المساهم الثاني} \\
\midrule
\textbf{الاسم} & --- & --- \\
\textbf{الصفة} & --- & --- \\
\textbf{التوقيع} & \rule{4cm}{0.4pt} & \rule{4cm}{0.4pt} \\
\textbf{التاريخ} & \rule{4cm}{0.4pt} & \rule{4cm}{0.4pt} \\
\bottomrule
\end{tabularx}
\end{table}

\end{document}
